\subsection{Factored Score $s_j(t) = \pi_j(t)\,u_j(t)$}
\label{sec:factored-score}

The linear combination in Eq.~\eqref{eq:score} conflates the two
orthogonal axes noted in Section~\ref{sec:recap} --- it is not clear, for
instance, why exploration should draw solely on historic data while
exploitation draws solely on current data. We instead factor the score
into a \emph{probability-of-opening} term and a \emph{utility-of-opening}
term, each independently schedulable:
\begin{equation}
\label{eq:pi-u}
\pi_j(t) = \alpha(t)\phi_j + (1-\alpha(t))f_j(t),
\qquad
u_j(t) = \beta(t)p_j + (1-\beta(t))\cdot 1,
\end{equation}
\begin{equation}
\label{eq:s-factored}
s_j(t) = \pi_j(t)\,u_j(t)
       = \Big(\alpha(t)\phi_j + [1-\alpha(t)]f_j(t)\Big)
         \Big(\beta(t)p_j + [1-\beta(t)]\Big).
\end{equation}
Here $\pi_j(t)$ blends historic and current opening-probability estimates,
while $u_j(t)$ blends the direct value of an open (the constant $1$) with
its indirect information value $p_j$. This factorization reuses every
primitive introduced above unchanged, so refinements from
Sections~\ref{sec:alpha-scheduling}--\ref{sec:variance-p} can be dropped
into $\pi_j(t)$ and $u_j(t)$ directly --- we use a confidence schedule
(Eq.~\eqref{eq:alpha-confidence}) for $\alpha(t)$ and a send-volume
schedule (Eq.~\eqref{eq:alpha-linear}/\eqref{eq:alpha-geometric}) for
$\beta(t)$.
